\section{Drone-Captured Self-Supervised Dataset}
\label{appendix:drone_dataset}

To support the self-supervised training stage of Brancher, a dedicated drone-based dataset was collected using a commercial \textit{Parrot Bebop 1} quadcopter.
\useImgWithCaption[]{Images/bebop.jpg}{0.25}{0.15}{The Parrot Bebop 1 drone used for data collection.} 
The goal of this acquisition was to obtain monocular aerial videos of human motion under realistic drone viewpoints, without relying on ground-truth 3D annotations.

The dataset consists of \textbf{698 video sequences}, each with a fixed temporal length of \textbf{81 frames}. All videos were recorded at a spatial resolution of $\mathbf{H=352}$ and $\mathbf{W=640}$ 
pixels. This fixed-length design ensures full compatibility with the temporal window used during training.

Data were collected across \textbf{five different locations} and involve \textbf{five distinct subjects}. The locations include varied outdoor environments, allowing changes in background, 
illumination, and scene geometry. The subjects were instructed to perform movements that generate highly diverse body configurations, covering a broad range of poses and limb articulations. This 
variability is crucial to avoid pose bias and to improve generalization under aerial viewpoints.

Importantly, each video focuses on a \textit{single person} captured from a drone perspective, resulting in monocular, single-view pose sequences. No 3D ground-truth annotations were recorded. The 
dataset was explicitly designed for \textbf{self-supervised training}, enabling the model to adapt to drone-specific viewpoints and motion dynamics without relying on supervised 3D labels.

This drone-captured dataset plays a central role in the second training stage of Brancher, where the network is optimized in a fully self-supervised manner to better generalize to real-world UAV 
scenarios.

\useImgWithCaption[]{Images/drone_dataset.jpg}{1.0}{0.6}{Some samples from the drone-captured dataset, showcasing the diversity of subjects, poses, and environments.}